\section{Introduction}
\label{sec:intro}

Steam methane reforming (SMR) still supplies most of the world's merchant and captive hydrogen, and it will remain the dominant route for at
least a decade while electrolytic capacity is built up \todo{support with a current IEA or similar figure and reference}. The core of an
SMR plant is a top-fired or side-fired furnace in which several hundred centrifugally cast high-alloy tubes, filled with nickel catalyst,
carry a mixture of natural gas and steam at 25--35 bar while their outer walls are heated to 1100--1200 K. The tubes operate in the creep
regime of the alloy by design, and their remaining life is the single most important maintenance variable of the plant: a tube
failure forces an unplanned shutdown, and a full retube costs of the order of \todo{quote a retube cost range with a reference}. Tube
life is in turn governed almost entirely by the outer-wall temperature at the hottest elevation, because the rupture time of the
alloy roughly halves for every 10--15 K of additional metal temperature \todo{confirm the rule of thumb for HP-Nb alloys and cite}.

Operating decisions therefore trade hydrogen production, fuel consumption and tube life against each other, but this trade-off is
rarely quantified. Plants are usually run against a maximum tube-wall temperature limit read from pyrometer surveys, and the effect of
a load change, a steam-to-carbon change or a burner adjustment on the tube life is left to experience. Two recent developments make a
quantitative approach both necessary and possible. On the demand side, hydrogen plants are increasingly expected to follow variable
downstream demand or to operate alongside intermittent renewable hydrogen, which raises the question of how much tube life
flexible operation actually costs. On the modelling side, validated first-principles models of the reformer tube
\cite{xu1989a,xu1989b,plehiers1989,pantoleontos2012} and of the furnace \cite{latham2011,zamaniyan2008} exist, machine-learning
surrogates make such models fast enough for hourly simulation and optimisation \cite{lao2016,tran2017,kumar2015,kumar2017}, and creep
data for the relevant alloys are available in public databases \cite{nims1990,nims1991}. What is missing is a coupled, calibrated and
uncertainty-aware chain from operating inputs to tube life, built entirely on published data and open code, so that its conclusions
can be checked and its data replaced.

\subsection{Previous work}

The reformer tube has been modelled at every level of complexity. Xu and Froment \cite{xu1989a} established the
Langmuir--Hinshelwood--Hougen--Watson kinetics of methane steam reforming and water-gas shift on a Ni/MgAl$_2$O$_4$ catalyst that
remains the standard intrinsic model, and in the companion paper \cite{xu1989b} they simulated an industrial tube with a
one-dimensional heterogeneous model and reported axial profiles of conversion, temperature and pressure for one operating point.
Plehiers and Froment \cite{plehiers1989} coupled the tube to a zone model of a top-fired furnace, and Pantoleontos et al.
\cite{pantoleontos2012} revisited the heterogeneous tube model with a dynamic formulation. Latham and co-workers
\cite{latham2008,latham2011} built a three-dimensional furnace model of a top-fired box reformer, calibrated it on plant measurements
including infrared tube-wall temperature surveys, and made the plant operating points and tube-wall temperatures available in the
thesis appendix \cite{latham2008}; that data set is the basis of the present calibration. Zamaniyan et al. \cite{zamaniyan2008}
proposed a simpler furnace model for control-oriented use.

On the materials side, the creep behaviour of the HP-Nb family of centrifugally cast tube alloys has been characterised in the
long-term data sheets of the National Institute for Materials Science \cite{nims1990,nims1991} and in remaining-life studies of
service-exposed tubes \cite{ray2003,whittaker2013,haidemenopoulos2021,guguloth2023}. Yeh \cite{yeh2021} combined a tube-wall
temperature estimate with a manufacturer's Larson--Miller curve for a proprietary HP-Nb alloy to estimate life, and the same
curve is used here as a placeholder until the public data sheets are transcribed. The Larson--Miller parameter \cite{larson1952}
and the linear life-fraction rule of Robinson \cite{robinson1952} remain the basis of design codes such as API 530 \cite{api530}.

Data-driven and hybrid approaches have been applied to reformers in the last decade: neural-network models of the reformer
\cite{lao2016,tran2017}, optimisation of reformer operation and furnace balancing \cite{kumar2015,kumar2017}, and, more recently,
physics-informed and digital-twin formulations \cite{yerimah2025,jafarizadeh2024,nabavi2025}. These studies generally address
production and efficiency; tube life, when it appears, is a constraint on a maximum temperature rather than an objective with a
quantified cost, and uncertainty in the kinetics, the heat transfer and the creep data is not propagated to the conclusions.

\subsection{Contribution and research questions}

This paper presents a physics-informed digital twin of an industrial top-fired SMR that links operating inputs to tube creep life
through a one-dimensional tube model, a long-furnace model, a Larson--Miller life module, machine-learning surrogates with
distribution-free prediction intervals, and Monte Carlo uncertainty propagation (Fig.~\ref{fig:schematic}). The tube model is verified
against the industrial case of Xu and Froment \cite{xu1989b}, the coupled model is calibrated on the plant data of Latham
\cite{latham2008} with cross-validation, and every number in the results is produced by a versioned open pipeline from the published
data. The twin is then used to answer four research questions:

\begin{enumerate}
\item[RQ1] Which operating variables govern the hot-spot temperature and the creep-life consumption of the tube, and how do they interact?
\item[RQ2] What is the tube-life cost, per unit of hydrogen produced, of flexible operation (load-following, renewable-following,
  over-firing events) and of catalyst ageing, relative to steady operation?
\item[RQ3] Which operating regimes are Pareto-optimal in hydrogen production, fuel consumption and life consumption, and how does the
  answer depend on how steam is charged?
\item[RQ4] How uncertain are these answers, which uncertainty sources dominate, and which conclusions survive the uncertainty?
\end{enumerate}

Section~\ref{sec:methods} describes the models, the data and the numerical methods. Section~\ref{sec:results} reports verification,
calibration and the answers to RQ1--RQ4. Section~\ref{sec:discussion} draws the industrial implications and states the limitations,
and Section~\ref{sec:conclusions} concludes.
